% Einheit 2 von 4 – Zeit (bank/einheiten/e2.jsonl, zusammenbau v0.1)
%% TODO Zweigzeile Teil 1: was hier gelernt wird (Satz in Schülersprache aus dem Einheitstitel) – Einheit 2
\einheitenkopf[e2]{Einheit 2 von 4 · Zeit}
\zweigzeile{ab Kl. 5 · P10}
\setcounter{aufgabe}{22}

%% TODO Titel: Ich-kann-Satz fehlt in der Bank (Platzhalter: Kettenname) – Nr. 23
%% TODO Anweisung: ein Satz über der ersten Teilaufgabe fehlt in der Bank – Nr. 23
\begin{aufgabe}{Zeit umrechnen und Zeitspannen berechnen}
%% TODO \rechenplatz{n}: Schrittzahl des Grundfalls fehlt in der Bank (nur bei fester Schreibform, 2.3 b)
\begin{teile}
\teil Der Film beginnt um 19:15 Uhr. Zeitpunkt oder Zeitspanne? Kreuze an.\\ \kreuz{Zeitpunkt}\\ \kreuz{Zeitspanne}
\teil $3$ min – wie viele s? \leerfeld[s]
\teil $240$ s – wie viele min? \leerfeld[min]
\teil Eine Viertelstunde – wie viele min? \leerfeld[min]
\teil $1{,}25$ h – wie viele min? \leerfeld[min]
\teil $90$ min – wie viele h? \leerfeld[h]
\teil Von 8:15 Uhr bis 9:40 Uhr – wie lange? \leerfeld[h] \leerfeld[min]
\teil Von 9:45 Uhr bis 12:15 Uhr – wie lange? \leerfeld[h] \leerfeld[min]
\teil Tim fährt um 8:30 Uhr los, fährt $1$ h $45$ min, macht $25$ min Pause und fährt dann noch $55$ min. Wann kommt er an? \leerfeld Uhr
\teil Eine Radtour dauert $5{,}5$ h. Wie viele Minuten sind das? \hfill (P10 2018 OS) \\ \leerfeld[min]
\teil Ein Zug fährt um 9:47 Uhr ab und ist $3$ h $25$ min unterwegs. Wann kommt er an? \hfill (P10 2021 OS) \\ \leerfeld Uhr
\end{teile}
\end{aufgabe}

%% TODO Titel: Ich-kann-Satz fehlt in der Bank (Platzhalter: Kettenname) – Nr. 24
%% TODO Anweisung: ein Satz über der ersten Teilaufgabe fehlt in der Bank – Nr. 24
\begin{aufgabe}{Zeiteinheiten nennen und ordnen}
\begin{teile}
\teil Ordne, mit der kleinsten beginnend: Tag, Sekunde, Stunde, Woche, Minute \leerfeld ; \leerfeld ; \leerfeld ; \leerfeld ; \leerfeld
\end{teile}
\end{aufgabe}

%% TODO Titel: Ich-kann-Satz fehlt in der Bank (Platzhalter: Kettenname) – Nr. 25
%% TODO Anweisung: ein Satz über der ersten Teilaufgabe fehlt in der Bank – Nr. 25
\begin{aufgabe}{Startzeit aus Endzeit und Dauer}
\begin{teile}
\teil Der Unterricht endet um 15:25 Uhr und dauert mit Pausen $6$ h $40$ min. Wann beginnt er? \leerfeld Uhr
\end{teile}
\end{aufgabe}

%% TODO Titel: Ich-kann-Satz fehlt in der Bank (Platzhalter: Kettenname) – Nr. 26
%% TODO Anweisung: ein Satz über der ersten Teilaufgabe fehlt in der Bank – Nr. 26
\begin{aufgabe}{Fahrplan: nächste Abfahrt, Fahrzeit}
\begin{teile}
\teil Mia ist um 7:25 Uhr an der Haltestelle Rathaus. Mit welchem Bus fährt sie, und wie lange braucht sie bis zur Schule?

\sachtabelle{lccc}{Halt & Bus 1 & Bus 2 & Bus 3}{Rathaus & 7:12 & 7:42 & 8:12\\ Bahnhof & 7:31 & 8:01 & 8:31\\ Schule & 7:48 & 8:18 & 8:48}
Bus \leerfeld , \leerfeld[min]
\end{teile}
\end{aufgabe}

%% TODO Titel: Ich-kann-Satz fehlt in der Bank (Platzhalter: Kettenname) – Nr. 27
%% TODO Anweisung: ein Satz über der ersten Teilaufgabe fehlt in der Bank – Nr. 27
\begin{aufgabe}{Dauer in eine sinnvolle Einheit bringen (Sekunden in Minuten, Stunden in Tage)}
\begin{teile}
\teil $72$ h – wie viele Tage? \leerfeld[Tage]
\end{teile}
\end{aufgabe}

%% TODO Titel: Ich-kann-Satz fehlt in der Bank (Platzhalter: Kettenname) – Nr. 28
%% TODO Anweisung: ein Satz über der ersten Teilaufgabe fehlt in der Bank – Nr. 28
\begin{aufgabe}{Zeit umrechnen und Zeitspannen berechnen – Fehler finden, Begründen, Anwendung, Darstellungswechsel}
\begin{teile}
\teil Ben schreibt: $2{,}75$ h $= 2$ h $75$ min. Wo ist der Fehler?
\teil Begründe: Warum sind $0{,}4$ h nicht $40$ min?
\teil Ein Handy-Akku lädt in $1{,}75$ h voll. Jonas steckt das Handy um 17:40 Uhr an und muss um 19:15 Uhr los. Ist der Akku dann voll?
\teil Ergänze die Tabelle für $1{,}75$ h.

\sachtabelle{ccc}{Dezimalstunde & Stunden und Minuten & Minuten}{1,75 h & \leerzelle & \leerzelle}
\end{teile}
\end{aufgabe}
